\chapter{基于锐度感知元经验回放的医学影像持续配准研究}
\label{chap:registration}
\setlength{\emergencystretch}{1em}

\section{引言}
\label{sec:registration-introduction}

前两章分别研究弱监督持续分割和类别扩展下的联邦持续分类。本章研究连续空间变换的持续学习：当器官、模态和配准关系依次变化时，模型在有限历史回放条件下学习新任务，并保持已见任务的配准能力。

医学图像配准根据固定图像和移动图像预测空间变换，使对应解剖位置对齐。本章按器官、模态组合和配准关系区分任务，各任务的预测对象均为空间变换，监督与损失配置分别说明。

直接在新配准任务上继续训练同一网络，容易使参数偏向当前图像外观和形变分布，降低早期任务的解剖重叠或地标对应精度。完整联合训练能够同时访问全部任务数据，但不符合阶段到达和历史访问受限条件。有限经验回放在两者之间保存少量历史固定--移动图像对，为旧任务提供直接梯度参照；其资源条件是允许保存和读取有限历史原始图像对。

为协调当前任务适应与历史任务保持，本章研究锐度感知元持续医学图像配准方法 SAMCL\cite{wang2024samcl}。SAMCL 在有限经验回放基础上，利用一阶元持续学习促进当前样本与历史样本更新方向的一致性，并把锐度感知优化嵌入元更新，使模型同时关注当前参数点及其邻域中的配准损失。实验按照 OASIS 脑部 MR--MR、腹部 CT--CT、NLST 肺部 CT--CT 和腹部 MR--CT 的固定顺序，比较直接顺序训练、参数正则、梯度投影、元经验回放以及完整数据访问参照，并分别报告 Dice、TRE、后向变化、回放容量和预设未见任务方向上的直接表现。

本章先给出任务增量配准的定义和 SAMCL 方法，再通过固定四任务序列分析配准精度、历史保持、回放容量与预设方向上的前向泛化。最后设计面向 MedCL 的模型或空间变换预测接入，明确图像对、坐标约定、评价参照和可视化之间的关系。

\section{任务增量持续医学图像配准的问题定义}
\label{sec:registration-formulation}

设第 $t$ 个配准任务的数据集为
$\mathcal{D}^{(t)}=\{(I_{F,i}^{(t)},I_{M,i}^{(t)})\}_{i=1}^{N_t}$，其中 $I_F$ 与 $I_M$ 分别表示固定图像和移动图像。配准网络
$g_{\theta}$ 根据图像对输出连续空间变换
$\phi_{\theta}^{(t)}=g_{\theta}(I_F^{(t)},I_M^{(t)})$，并通过空间重采样得到变形后的移动图像 $I_M\circ\phi_{\theta}^{(t)}$。不同任务可以具有不同器官、模态、个体间或个体内关系、图像外观、形变分布和监督信息，因此其配准损失记为
$\mathcal{L}^{(t)}(\phi_{\theta};I_F,I_M)$。

若所有已见任务数据可以一并访问，则完成前 $K$ 个任务后的理想累积目标为
\begin{equation}
\widehat{\theta}^{(K)}
=
\operatorname*{arg\,min}_{\theta}
\sum_{t=1}^{K}
\mathbb{E}_{(I_F,I_M)\sim\mathcal{D}^{(t)}}
\left[
\mathcal{L}^{(t)}(\phi_{\theta};I_F,I_M)
\right],
\qquad K=1,\ldots,T.
\label{eq:samcl-cumulative-objective}
\end{equation}
在顺序训练条件下，当前阶段通常只能完整访问 $\mathcal{D}^{(K)}$。若直接在新任务上继续优化，参数可能偏向当前分布和配准目标，并降低旧任务上的解剖重叠或地标对应精度。

四项任务均属于配准，但器官、模态组合、配准关系和评价对象不同，形成连续的分析目标变化。实验使用一个顺序更新的配准模型，并在各任务自身的图像对和评价参照上检验其能力。

为提供有限历史信息，SAMCL 使用容量受限的历史回放集合 $\mathcal{M}$ 保存少量历史固定--移动图像对，满足
$|\mathcal{M}|\leq b$。这一资源条件介于联合训练与无历史原始样本回放之间：模型不能重新调用完整历史训练集，但后续更新会读取 $\mathcal{M}$ 中保留的历史原始图像对。有限回放用于提供旧任务梯度，元持续学习用于减小当前与历史梯度之间的干扰，锐度感知优化则约束参数邻域中的损失。模型的历史保持和前向泛化只能在给定任务顺序、回放容量、训练预算及评价数据上验证。

\section{锐度感知元持续医学图像配准方法}
\label{sec:samcl}

\subsection{跨器官与跨模态任务变化中的配准更新与知识保持}
\label{subsec:samcl-motivation}

医学图像配准通过估计固定图像与移动图像之间的空间变换，使两幅图像中的对应解剖结构在同一坐标系下对齐，是疾病进展监测、术前规划、手术导航和多模态影像联合分析中的基础环节\cite{maintz1998registration,gupta2021registration,maes1997mutualinformation,sotiras2013deformable}。深度学习方法能够直接从图像对学习配准变换，在推理效率和任务内配准精度方面取得了较好结果\cite{haskins2020deepreg,balakrishnan2019voxelmorph}。然而，现有网络通常围绕特定解剖区域、成像模态和配准关系构建，例如脑部 MR--MR、肺部 CT--CT 或腹部 CT--CT 配准\cite{kang2022dual,che2023amnet,hering2021lungreg}。一旦测试数据的器官、外观、模态关系或采集条件偏离训练分布，任务特定模型的性能容易下降。

真实临床数据还具有持续到达的特点。新设备投入使用、检查方案变化以及新的器官或模态进入分析流程，都会使配准任务随时间扩展。若为每个任务分别保存一套网络，模型数量和维护成本会随任务增加；若等待全部数据集中后再进行多任务训练，又要求不同阶段的数据能够被长期保存并集中访问，这在医疗数据治理和隐私要求下往往难以实现\cite{lee2020clinicalcl,vokinger2021medicaldevices}。直接使用当前任务继续更新同一网络虽然满足顺序训练条件，却会因分布变化破坏早期任务的参数解，产生灾难性遗忘\cite{thrun1995lifelong,mccloskey1989catastrophic,delange2022continual}。

本文将不同器官、模态和配准关系的数据集组织为顺序任务，在不能集中访问全部历史数据的条件下持续更新单一模型，并评价已见任务保持及预设方向上的泛化表现\cite{wang2024samcl}。

图\ref{fig:samcl-framework}展示了这一设定。训练阶段依次到达脑部 MR--MR、腹部 CT--CT、肺部 CT--CT 和腹部 MR--CT 图像对，任务之间同时存在解剖区域、图像外观、模态关系和个体内/个体间配准方式的变化。推理阶段不再为各任务维护独立网络，而是使用顺序训练后得到的单一模型处理已见任务。SAMCL 以有限经验回放保存少量历史图像对，以一阶元持续学习协调当前和历史样本的更新方向，并进一步通过锐度感知优化改善参数邻域内的损失，从而同时面向历史保持和前向泛化。

\IfFileExists{figures/ch06/samcl_framework.pdf}{%
\begin{figure}[htbp]
\centering
\includegraphics[width=\textwidth]{ch06/samcl_framework.pdf}
\caption[SAMCL 的持续医学图像配准设定]{SAMCL 的持续医学图像配准设定。模型依次学习脑部 MR--MR、腹部 CT--CT、肺部 CT--CT 和腹部 MR--CT 配准，完成顺序训练后使用同一模型处理各已见任务。}
\label{fig:samcl-framework}
\end{figure}
}{}

\subsection{配准网络与持续学习目标}
\label{subsec:samcl-registration-objective}

设固定图像为 $f:\Omega_F\rightarrow\mathbb{R}$，移动图像为 $m:\Omega_M\rightarrow\mathbb{R}$。医学图像配准需要估计从固定图像域到移动图像域的空间变换
$\phi:\Omega_F\rightarrow\Omega_M$，使变形后的移动图像
$m\circ\phi$ 与固定图像中的对应结构对齐。学习式配准网络以图像对为输入，并用参数 $\boldsymbol{\theta}$ 表示该变换：
\begin{equation}
\phi_{\boldsymbol{\theta}}
=
g(f,m;\boldsymbol{\theta}).
\label{eq:samcl-registration-network}
\end{equation}
对单一数据集 $\mathcal{D}=\{(f_i,m_i)\}_{i=1}^{N}$，模型通过最小化经验风险学习参数：
\begin{equation}
\widehat{\boldsymbol{\theta}}
=
\operatorname*{arg\,min}_{\boldsymbol{\theta}}
\mathbb{E}_{(f,m)\sim\mathcal{D}}
\left[
\mathcal{L}(\phi_{\boldsymbol{\theta}};f,m)
\right],
\label{eq:samcl-single-task-objective}
\end{equation}
其中，配准损失可统一写为
\begin{equation}
\mathcal{L}(\phi_{\boldsymbol{\theta}};f,m)
=
D\!\left(f,m\circ\phi_{\boldsymbol{\theta}}\right)
+
R\!\left(\phi_{\boldsymbol{\theta}}\right).
\label{eq:samcl-registration-loss}
\end{equation}
$D(\cdot,\cdot)$ 衡量固定图像与变形后移动图像的外观或解剖不一致，$R(\cdot)$ 用于约束形变的平滑性。不同任务能够提供的标签或地标不同，因此第 $t$ 个任务使用任务特定损失 $\mathcal{L}^{(t)}$。

在持续配准中，数据集
$\{\mathcal{D}^{(t)}\}_{t=1}^{T}$ 按顺序到达。若完成前 $K$ 个任务后仍能访问所有已见数据，则理想目标为式\eqref{eq:samcl-cumulative-objective}中的累积风险最小化。实际训练阶段通常只能完整访问当前数据集 $\mathcal{D}^{(K)}$，因此需要在有限历史访问条件下近似这一目标。SAMCL 保存容量受限的历史回放集合 $\mathcal{M}$，后续阶段能够重新读取其中的固定--移动图像对，但不能访问完整历史训练集。该设定介于集中式多任务训练和无历史原始样本回放持续学习之间。

\subsection{有限经验回放下的元持续学习}
\label{subsec:samcl-meta-continual}

仅将当前任务损失与回放损失相加，不能保证两类更新方向相互兼容。当当前任务梯度与历史任务梯度方向相反时，一次有利于新任务的参数更新可能增大旧任务损失。元持续学习从过去知识保持和未来任务适应两个方向组织更新，通过一阶元学习鼓励不同样本或任务之间的迁移，并减小梯度干扰。SAMCL 以 Meta Experience Replay（MER）为基础，MER 结合经验回放与 Reptile 一阶元学习\cite{riemer2019mer,nichol2018firstorder,method_replay_er}。

具体地，令 $\mathcal{Q}_t$ 表示由当前任务数据与历史回放集合共同形成的采样分布，并从中得到样本序列
$z_i=(f_i,m_i)$。样本对应任务记为 $t_i$，内层学习率为 $\alpha$。MER 近似优化下式：
\begin{equation}
\begin{aligned}
\min_{\boldsymbol{\theta}}
\mathbb{E}_{z_1,\ldots,z_s\sim\mathcal{Q}_t}
\Bigg[
\sum_{i=1}^{s}
\Bigg(&
\mathcal{L}^{(t_i)}_i(\boldsymbol{\theta})
-
\frac{\alpha}{2}
\sum_{j=1}^{i-1}
\nabla_{\boldsymbol{\theta}}
\mathcal{L}^{(t_i)}_i(\boldsymbol{\theta})^{\top}
\nabla_{\boldsymbol{\theta}}
\mathcal{L}^{(t_j)}_j(\boldsymbol{\theta})
\Bigg)
\Bigg].
\end{aligned}
\label{eq:samcl-meta-objective}
\end{equation}
第一项降低各样本自身的配准损失；第二项显式体现梯度内积。若两个梯度的内积为正，沿其中一个样本的下降方向更新通常也有利于另一样本；若内积为负，则表明两次更新存在干扰。因此，一阶元更新不只追求当前批次损失下降，还通过训练轨迹近似增加跨任务梯度的一致性，使模型在有限记忆下兼顾新任务学习和历史任务保持。

为在固定容量下保留历史样本，SAMCL 使用水库采样更新历史回放集合。设到当前时刻已经观察到 $n$ 个图像对，回放容量为 $B_M$。当回放集合未满时，当前图像对直接写入；回放集合填满后，第 $n$ 个样本以 $B_M/n$ 的概率进入回放集合，并随机替换其中一个已有样本。由此，顺序数据流中的样本具有近似相同的保留机会，回放集合不需要预先知道任务长度。水库采样只控制历史样本的进入方式，并不能保证不同器官、模态或任务在有限容量下完全均衡，因此元更新的效果仍依赖回放集合对历史分布的代表性。

\subsection{锐度感知元持续学习}
\label{subsec:samcl-sharpness-update}

可变形配准损失通常高度非凸，可能包含多个局部极小值\cite{hill2001registration}。在有限样本和顺序分布变化下，仅在当前参数点取得较低损失，并不意味着其邻域内的损失同样稳定；如果参数附近的损失曲面较尖锐，轻微的任务变化、批次采样差异或后续更新都可能导致性能明显下降。为此，SAMCL 将锐度感知最小化引入元持续学习\cite{foret2021sam}。

对任务 $t$，锐度感知损失定义为
\begin{equation}
\mathcal{L}^{(t)}_{\mathrm{sam}}(\boldsymbol{\theta})
=
\max_{\|\boldsymbol{\epsilon}\|_2\leq\rho}
\mathcal{L}^{(t)}
\left(\phi_{\boldsymbol{\theta}+\boldsymbol{\epsilon}}\right),
\label{eq:samcl-sharpness-loss}
\end{equation}
其中，$\rho$ 为参数邻域半径，$\boldsymbol{\epsilon}$ 为参数扰动。该目标最小化半径 $\rho$ 邻域内的最坏损失，实际优化通过局部近似完成，以寻找邻域损失较低的参数区域。

将式\eqref{eq:samcl-sharpness-loss}代入元持续学习目标，可得 SAMCL 的总体目标：
\begin{equation}
\begin{aligned}
\min_{\boldsymbol{\theta}}
\mathbb{E}_{z_1,\ldots,z_s\sim\mathcal{Q}_t}
\Bigg[
\sum_{i=1}^{s}
\Bigg(&
\mathcal{L}^{(t_i)}_{\mathrm{sam},i}(\boldsymbol{\theta})
-
\frac{\alpha}{2}
\sum_{j=1}^{i-1}
\nabla_{\boldsymbol{\theta}}
\mathcal{L}^{(t_i)}_{\mathrm{sam},i}(\boldsymbol{\theta})^{\top}
\nabla_{\boldsymbol{\theta}}
\mathcal{L}^{(t_j)}_{\mathrm{sam},j}(\boldsymbol{\theta})
\Bigg)
\Bigg].
\end{aligned}
\label{eq:samcl-full-objective}
\end{equation}
该式同时包含两个层面的约束：梯度内积项面向当前样本与历史样本之间的迁移和干扰，锐度感知损失则面向每个样本对应的局部损失几何。二者分别作用于跨样本更新关系和参数邻域稳定性。

内层最大化通常难以精确求解，因此采用局部一阶近似。对
$\mathcal{L}^{(t)}(\boldsymbol{\theta}+\boldsymbol{\epsilon})$
在 $\boldsymbol{\epsilon}=\boldsymbol{0}$ 处作一阶泰勒展开：
\begin{equation}
\mathcal{L}^{(t)}(\boldsymbol{\theta}+\boldsymbol{\epsilon})
\approx
\mathcal{L}^{(t)}(\boldsymbol{\theta})
+
\boldsymbol{\epsilon}^{\top}
\nabla_{\boldsymbol{\theta}}
\mathcal{L}^{(t)}(\boldsymbol{\theta}).
\label{eq:samcl-taylor}
\end{equation}
在 $\ell_2$ 邻域约束下，上式的最大化可写为对偶范数问题，其近似最优扰动为
\begin{equation}
\widehat{\boldsymbol{\epsilon}}(\boldsymbol{\theta})
=
\rho
\frac{
\nabla_{\boldsymbol{\theta}}
\mathcal{L}^{(t)}(\boldsymbol{\theta})
}{
\left\|
\nabla_{\boldsymbol{\theta}}
\mathcal{L}^{(t)}(\boldsymbol{\theta})
\right\|_2
}.
\label{eq:samcl-perturbation}
\end{equation}
实际更新先在当前参数处计算梯度并构造
$\widehat{\boldsymbol{\epsilon}}$，再在扰动参数
$\boldsymbol{\theta}+\widehat{\boldsymbol{\epsilon}}$
处计算用于下降的梯度，即
$\nabla_{\boldsymbol{\theta}}\mathcal{L}^{(t)}_{\mathrm{sam}}$
由扰动点的梯度进行近似。SAMCL 依次对当前批次和回放批次执行锐度感知更新，再从更新前参数向内层更新后的参数插值，完成一阶元更新。有限回放保留旧任务信息，元更新协调当前与回放数据的学习，锐度感知目标约束参数扰动下的损失。

\subsection{训练流程与配准模型}
\label{subsec:samcl-algorithm}

算法\ref{alg:samcl}给出 SAMCL 的训练流程。第 $t$ 个数据集到达后，模型从上一任务参数
$\widehat{\boldsymbol{\theta}}^{(t-1)}$ 继续训练。对当前图像对，算法首先保存更新前参数 $\boldsymbol{\theta}_0$，随后根据式\eqref{eq:samcl-perturbation}执行当前任务的锐度感知内层更新；当前图像对通过水库采样写入历史回放集合。然后，从回放集合采样图像对，按照其所属任务损失再次执行锐度感知更新。最后使用元学习率 $\beta$，从更新前参数 $\boldsymbol{\theta}_0$ 向内层更新后的参数 $\boldsymbol{\theta}$ 插值，完成一阶元更新。任务结束时保存
$\widehat{\boldsymbol{\theta}}^{(t)}$，并作为下一任务的初始化。

\begin{algorithm}[htbp]
\caption{锐度感知元持续医学图像配准（SAMCL）}
\label{alg:samcl}
\KwData{顺序到达的训练集 $\mathcal{D}^{(t)}=\{(f_i,m_i)\}_{i=1}^{N_t}$，$t=1,\ldots,T$}
\KwIn{内层学习率 $\alpha$，元学习率 $\beta$，邻域半径 $\rho$，内层批大小 $s$}
\KwOut{完成全部任务后的参数 $\widehat{\boldsymbol{\theta}}^{(T)}$}
随机初始化 $\widehat{\boldsymbol{\theta}}^{(0)}$，令历史回放集合 $\mathcal{M}\leftarrow\varnothing$\;
\For{$t=1,\ldots,T$}{
  将 $\mathcal{D}^{(t)}$ 组织为当前任务批次 $\mathcal{B}_t=\{B_i\}_{i=1}^{b_t}$，令 $\boldsymbol{\theta}\leftarrow\widehat{\boldsymbol{\theta}}^{(t-1)}$\;
  \For{$i=1,\ldots,b_t$}{
    $\boldsymbol{\theta}_0\leftarrow\boldsymbol{\theta}$\;
    \For{$j=1,\ldots,s$}{
      $\boldsymbol{g}\leftarrow\nabla_{\boldsymbol{\theta}}\mathcal{L}^{(t)}(\boldsymbol{\theta};B_i[j])$，$\widehat{\boldsymbol{\epsilon}}\leftarrow\rho\boldsymbol{g}/\|\boldsymbol{g}\|_2$\;
      $\boldsymbol{\theta}\leftarrow\boldsymbol{\theta}-\alpha\nabla_{\boldsymbol{\theta}}\mathcal{L}^{(t)}(\boldsymbol{\theta}+\widehat{\boldsymbol{\epsilon}};B_i[j])$\;
    }
    使用水库采样将当前图像对写入 $\mathcal{M}$\;
    从 $\mathcal{M}$ 中采样回放批次 $B_i^{\mathcal{M}}$，其所属任务记为 $t_i$\;
    $\boldsymbol{g}_{\mathcal{M}}\leftarrow\nabla_{\boldsymbol{\theta}}\mathcal{L}^{(t_i)}(\boldsymbol{\theta};B_i^{\mathcal{M}})$，$\widehat{\boldsymbol{\epsilon}}_{\mathcal{M}}\leftarrow\rho\boldsymbol{g}_{\mathcal{M}}/\|\boldsymbol{g}_{\mathcal{M}}\|_2$\;
    $\boldsymbol{\theta}\leftarrow\boldsymbol{\theta}-\alpha\nabla_{\boldsymbol{\theta}}\mathcal{L}^{(t_i)}(\boldsymbol{\theta}+\widehat{\boldsymbol{\epsilon}}_{\mathcal{M}};B_i^{\mathcal{M}})$\;
    $\boldsymbol{\theta}\leftarrow\boldsymbol{\theta}_0+\beta(\boldsymbol{\theta}-\boldsymbol{\theta}_0)$\tcp*{一阶元更新}
  }
  $\widehat{\boldsymbol{\theta}}^{(t)}\leftarrow\boldsymbol{\theta}$\;
}
\Return{$\widehat{\boldsymbol{\theta}}^{(T)}$}\;
\end{algorithm}

配准骨干采用 Large Kernel U-Net（LK-UNet），首层包含 16 个卷积核\cite{jia2022lkunet}。网络输出平稳速度场，并通过微分同胚积分获得空间变换；固定图像和移动图像之间采用对称配准。形变正则项由膜能和弯曲能组成，用于约束局部位移变化和形变曲率\cite{ashburner2007diffeomorphic}。这些正则项控制形变平滑性，而 SAMCL 的元更新和锐度感知目标分别面向跨任务干扰与损失曲面，两者不能替代配准模型自身的形变约束。

\section{实验设计与结果分析}
\label{sec:registration-experiments}

\subsection{实验材料与任务序列}
\label{subsubsec:samcl-materials}

实验采用 Learn2Reg 挑战中的四项三维医学图像配准任务，并按照 OASIS 脑部 MR--MR、腹部 CT--CT、NLST 肺部 CT--CT 和腹部 MR--CT 的顺序到达\cite{hering2021learn2reg}。OASIS 来源于开放获取脑部 MRI 数据集，用于脑部个体间配准\cite{marcus2007oasis}；腹部 CT--CT 任务用于腹部个体间配准\cite{xu2016abdomenreg}；NLST 肺部 CT--CT 与腹部 MR--CT 数据来自公开影像资源，用于个体内肺部配准和跨模态腹部配准\cite{clark2013tcia,xu2016abdomenreg}。任务序列同时改变了解剖区域、图像外观、成像模态以及个体内/个体间关系，从而构成较强的连续分布变化。

相应地，表\ref{tab:samcl-datasets}给出四项任务的数据组织。OASIS、腹部 CT--CT 和腹部 MR--CT 分别使用 35、13 和 9 个对应解剖标签评价配准后的结构重叠；NLST 使用对应地标评价目标配准误差。表中的训练、验证和测试数量与本文前期实验设置保持一致\cite{wang2024samcl}，数量按原论文的三维影像数报告。

\begin{table}[htbp]
\centering
\caption{SAMCL 四个顺序配准任务的数据组织、预处理与评价。训练、验证和测试数量按原论文的影像数报告。}
\label{tab:samcl-datasets}
\thesistablefont
\renewcommand{\arraystretch}{1.14}
\setlength{\tabcolsep}{3pt}
\begin{tabularx}{\textwidth}{@{}cXp{1.8cm}p{2.35cm}p{1.65cm}p{2.0cm}p{1.5cm}@{}}
\toprule
顺序 & 数据集与任务 & 配准关系 & 训练/验证/测试 & 重采样间距 & 标签或地标 & 评价指标 \\
\midrule
1 & OASIS，脑部 MR--MR & 个体间 & 290/40/84 & $2\times2\times2$ mm & 35 类解剖标签 & 平均 Dice \\
2 & 腹部 CT--CT & 个体间 & 21/3/6 & $3\times3\times3$ mm & 13 类解剖标签 & 平均 Dice \\
3 & NLST，肺部 CT--CT & 个体内 & 294/42/84 & $3\times3\times3$ mm & 对应地标 & TRE（mm） \\
4 & 腹部 MR--CT & 个体内、跨模态 & 80/9/16 & $3\times3\times3$ mm & 9 类解剖标签 & 平均 Dice \\
\bottomrule
\end{tabularx}
\end{table}

所有体数据均裁剪或填充为
$112\times96\times112$。OASIS 重采样到
$2\times2\times2$ mm，其余任务重采样到
$3\times3\times3$ mm。CT 图像强度裁剪到
$[-200,300]$；随后，CT 和 MR 图像均线性归一化到 $[0,1]$。为避免不同方法因数据增强策略不同而获得额外收益，训练阶段不使用数据增强。

\subsection{对比方法与历史访问条件}
\label{subsubsec:samcl-baselines}

实验从任务特定训练、集中式训练、直接顺序训练和持续学习四个层面设置对比方法。未配准输入用于给出固定图像与移动图像在变形前的基础差异；独立训练和集中式多任务训练分别作为任务特定参考与完整数据访问参考；直接顺序训练用于观察不采用抗遗忘机制时的性能变化。持续学习方法包括 EWC、SI、GPM 和 MER。

EWC 根据旧任务估计参数重要性，并通过 Fisher 信息加权的二次惩罚限制重要参数偏移\cite{kirkpatrick2017ewc}。SI 在训练过程中累计参数对历史损失下降的贡献，以在线重要性约束后续更新\cite{method_regu_si}。GPM 从历史任务的表示中构建重要子空间，将新任务梯度投影到其正交补，以减小对旧任务方向的干扰\cite{saha2021gpm}。上述三种方法不回放历史原始图像。MER 使用有限经验回放和一阶元学习，通过增加跨样本梯度对齐来兼顾迁移与遗忘\cite{riemer2019mer}。SAMCL 与 MER 使用相同类型的有限历史图像对访问条件，其差别在于 SAMCL 将锐度感知更新嵌入当前和回放轨迹。

表\ref{tab:samcl-baselines}汇总各方法的训练组织和历史数据访问方式。独立训练为每项任务分别训练模型，集中式多任务训练同时使用全部任务数据；两者分别作为任务特定训练和完整数据访问条件下的性能参照，不与受限持续学习方法视为相同训练条件。

\begin{table}[htbp]
\centering
\caption{SAMCL 实验中的对比方法及其训练信息访问条件}
\label{tab:samcl-baselines}
\thesistablefont
\renewcommand{\arraystretch}{1.16}
\setlength{\tabcolsep}{3pt}
\begin{tabular}{p{1.8cm}p{4.6cm}p{3.2cm}p{2.8cm}}
\toprule
方法 & 训练组织 & 历史原始图像访问 & 模型组织 \\
\midrule
未配准 & 不估计空间变换，直接评价固定图像与移动图像 & 不涉及 & 不涉及 \\
独立训练 & 每个任务仅使用本任务数据单独训练 & 不跨任务访问 & 每任务一个模型 \\
集中式多任务 & 所有任务数据同时参与联合训练 & 完整访问全部任务数据 & 单一模型 \\
直接顺序训练 & 当前任务到达后直接继续优化 & 不访问历史图像 & 单一持续模型 \\
EWC & 以参数重要性正则限制重要参数变化 & 不访问历史图像 & 单一持续模型 \\
SI & 以在线累计的重要性约束参数变化 & 不访问历史图像 & 单一持续模型 \\
GPM & 通过历史重要子空间投影新任务梯度 & 不访问历史图像 & 单一持续模型 \\
MER & 有限经验回放与一阶元持续学习 & 容量受限的历史图像对 & 单一持续模型 \\
SAMCL & 有限经验回放、元持续学习与锐度感知更新 & 容量受限的历史图像对 & 单一持续模型 \\
\bottomrule
\end{tabular}
\end{table}

\subsection{评价指标与实现细节}
\label{subsubsec:samcl-implementation}

对于 OASIS、腹部 CT--CT 和腹部 MR--CT，实验在相应解剖标签上计算平均 Dice；对于 NLST，在对应地标上计算以毫米为单位的目标配准误差（target registration error, TRE）。令 $R_{i,j}$ 表示完成任务 $i$ 后在任务 $j$ 测试集上的评价，则任务 $j$ 在全部任务训练结束后的最终性能和后向变化定义如下。其中，$\mathrm{AVG}_j$ 沿用原研究记号，表示任务 $j$ 内部的平均性能 $R_{T,j}$，不是四个任务指标的平均值：
\begin{equation}
\mathrm{AVG}_j=R_{T,j},
\qquad
\mathrm{BWT}_j=R_{T,j}-R_{j,j}.
\label{eq:samcl-avg-bwt}
\end{equation}
对 Dice，BWT 为负表示重叠精度下降，为正表示提高，因此 Dice-BWT 越高越好。对 TRE，BWT 为正表示误差增大，为负表示误差减小，因此 TRE-BWT 越低越好。最后一个任务之后没有进一步训练，其 BWT 按定义为零，不能据此判断方法的历史保持能力。Dice 与 TRE 的单位不同，本节分别报告各任务结果。

在实现方面，所有方法使用相同的 LK-UNet 骨干。实验在 PyTorch 中实现，并在 NVIDIA RTX 3090 GPU 上运行。每个任务训练 $10000$ 次迭代，采用 Adam 优化器\cite{kingma2015adam}，学习率为 $1\times10^{-4}$，批大小为 $4$。SAMCL 的历史回放容量设为 $200$，元学习率 $\beta$ 设为 $0.25$，锐度邻域半径 $\rho$ 设为 $0.05$。

四项任务的标注形式不同，因此训练损失随任务变化。OASIS、腹部 CT--CT 和腹部 MR--CT 的损失由窗口大小为 $3$ 的负局部归一化互相关（local normalized cross-correlation, LNCC）和软 Dice 损失组成\cite{avants2008syn,hu2018weakreg}；NLST 的损失由负 LNCC 和基于对应地标的距离损失组成。所有任务均加入膜能和弯曲能构成的形变正则项，两项正则权重均设为 $1$。配准在固定图像与移动图像之间对称执行，并使用由平稳速度场参数化的微分同胚变换。

\subsection{配准精度与知识保持分析}
\label{subsubsec:samcl-analysis}

表\ref{tab:samcl-main-results}给出四任务顺序训练结束后的最终性能和后向变化。直接顺序训练在前三个任务上出现明显知识损失：OASIS 和腹部 CT--CT 的 Dice-BWT 分别为 $-0.395$ 和 $-0.129$，NLST 的 TRE-BWT 为 $3.508$ mm。SAMCL 将三项后向变化改善为 $-0.037$、$-0.039$ 和 $0.306$ mm，说明有限回放、元持续更新与锐度感知训练的组合能够明显缓解直接顺序训练的遗忘。

\begin{table}[htbp]
\centering
\caption{SAMCL 四任务顺序实验的最终性能与后向变化。}
\label{tab:samcl-main-results}
\thesistablefont
\renewcommand{\arraystretch}{1.12}
\setlength{\tabcolsep}{3pt}
\begin{tabular}{@{}lcccc@{}}
\toprule
\multicolumn{5}{c}{（a）最终性能 $R_{T,j}$} \\
方法 & OASIS Dice & 腹部 CT Dice & NLST TRE & 腹部 MR--CT Dice \\
\midrule
未配准 & $\underline{0.598}^{*}$ & $\underline{0.311}^{*}$ & $\underline{7.932}^{*}$ & $\underline{0.353}^{*}$ \\
独立训练 & $0.738^{*}$ & $\underline{0.335}^{*}$ & $3.259^{*}$ & $\underline{0.550}$ \\
集中式多任务 & $0.743^{*}$ & 0.358 & $2.662^{*}$ & $\underline{0.506}$ \\
直接顺序训练 & $\underline{0.348}^{*}$ & $\underline{0.270}$ & $\underline{6.922}^{*}$ & 0.568 \\
EWC & $\underline{0.659}^{*}$ & $\best{0.364^{*}}$ & 3.836 & $\underline{0.506}$ \\
SI & $\underline{0.461}^{*}$ & $\underline{0.303}^{*}$ & $\underline{8.186}^{*}$ & $\underline{0.555}^{*}$ \\
GPM & $\underline{0.458}^{*}$ & $\underline{0.290}^{*}$ & $\underline{7.095}^{*}$ & $\underline{0.556}^{*}$ \\
MER & \best{0.698} & $\underline{0.331}^{*}$ & $\best{3.611^{*}}$ & $\underline{0.532}$ \\
SAMCL & 0.697 & 0.352 & 3.855 & \best{0.563} \\
\bottomrule
\end{tabular}

\vspace{8pt}
\begin{tabular}{@{}lcccc@{}}
\toprule
\multicolumn{5}{c}{（b）后向变化 $R_{T,j}-R_{j,j}$} \\
方法 & OASIS Dice & 腹部 CT Dice & NLST TRE & 腹部 MR--CT Dice \\
\midrule
直接顺序训练 & $\underline{-0.395}^{*}$ & $\underline{-0.129}^{*}$ & $\underline{3.508}^{*}$ & 0  \\
EWC & $\underline{-0.083}^{*}$ & $\best{-0.011^{*}}$ & $\underline{0.386}$ & 0 \\
SI & $\underline{-0.281}^{*}$ & $\underline{-0.084}^{*}$ & $\underline{4.685}^{*}$ & 0 \\
GPM & $\underline{-0.281}^{*}$ & $\underline{-0.085}^{*}$ & $\underline{3.574}^{*}$ & 0 \\
MER & $\best{-0.036}$ & $\underline{-0.040}$ & $\best{0.077^{*}}$ & 0 \\
SAMCL & $-0.037$ & $-0.039$ & 0.306 & 0 \\
\bottomrule
\end{tabular}

\vspace{3pt}
\begin{minipage}{0.98\textwidth}
\thesistablefont
注：Dice 越高越好，NLST TRE（mm）越低越好。红色粗体标出 EWC、SI、GPM、MER 和 SAMCL 中的列最优值；下划线表示该值劣于 SAMCL，星号表示与 SAMCL 的配对 $t$ 检验达到 $p<0.05$。未配准、独立训练和集中式多任务不定义后向变化；最后任务的零后向变化不作最优标记。
\end{minipage}
\end{table}

从最终配准精度看，SAMCL 在 OASIS、腹部 CT--CT 和腹部 MR--CT 上分别取得 $0.697$、$0.352$ 和 $0.563$ 的 Dice，在 NLST 上取得 $3.855$ mm 的 TRE。相较 SI 和 GPM，SAMCL 在 OASIS 与腹部 CT--CT 上具有更高 Dice，并明显减小早期任务的后向下降。与 EWC 相比，SAMCL 在 OASIS 上具有更高最终 Dice 和更小后向下降；腹部 CT--CT 上 EWC 的最终 Dice 略高，且后向下降更小。该结果说明参数正则化在某一任务上仍可能有效，但其表现会随器官和损失形式变化。

MER 与 SAMCL 均使用有限回放。两者在 OASIS 最终 Dice 和前两个任务 BWT 上接近；SAMCL 在腹部 CT--CT 和腹部 MR--CT 上取得更高 Dice，MER 在 NLST 上取得更低 TRE 和更小 TRE-BWT。锐度感知更新的收益随配准任务和评价指标变化。

独立训练和集中式多任务训练可以分别访问单一任务完整数据或全部任务数据。SAMCL 在 OASIS 上低于两项参考，在腹部 CT--CT 上接近集中式多任务训练，并在腹部 MR--CT 上高于独立训练和集中式多任务训练。直接顺序训练在最后任务上取得 $0.568$ 的 Dice，略高于 SAMCL 的 $0.563$；由于最后一个任务之后没有进一步训练，该差异反映当前任务适应，而不能用于评价历史保持。在前三个历史任务上，SAMCL 缩小了直接顺序训练与完整数据访问参考之间的差距；有限历史回放仍不能替代集中式访问全部历史任务。

图\ref{fig:samcl-qualitative}给出 OASIS 与腹部 CT--CT 的代表性图像对。每个任务的上排显示固定图像与配准后移动图像之间的差异图，下排显示不同方法得到的配准结果及对应 Dice。直接顺序训练、SI 和 GPM 在早期任务上出现较明显差异，而 SAMCL 的差异图和变形后图像更接近独立训练与集中式多任务训练。该可视化与表\ref{tab:samcl-main-results}中的早期任务保持结果一致。

\IfFileExists{figures/ch06/samcl_qualitative.pdf}{%
\begin{figure}[htbp]
\centering
\includegraphics[width=\textwidth]{ch06/samcl_qualitative.pdf}
\caption[不同训练策略的代表性配准结果]{不同训练策略在 OASIS 与腹部 CT--CT 代表性图像对上的配准结果。每个任务左侧给出固定图像与移动图像；各方法的上排为固定图像与配准后移动图像之间的差异图，下排为配准结果及对应 Dice。CT 图像显示时将强度裁剪到 $[-200,300]$。}
\label{fig:samcl-qualitative}
\end{figure}
}{}

\subsection{回放容量与域内、跨任务泛化分析}
\label{subsubsec:samcl-ablation-generalization}

图\ref{fig:samcl-ablation-generalization}(a)比较回放容量为 $50$、$100$ 和 $200$ 时的配准结果。OASIS、腹部 CT--CT 和腹部 MR--CT 三项 Dice 任务的最终精度总体随容量增加而提高。该比较说明回放容量会影响最终性能，方法比较时需同时报告保存的图像对数量。

图\ref{fig:samcl-ablation-generalization}(b)比较采用与不采用锐度感知最小化时的泛化表现。CTCT$\rightarrow$CTCT 和 MRCT$\rightarrow$MRCT 在同一配准任务的未参与训练数据上测试；CTCT$\rightarrow$MRCT、CTCT$\rightarrow$NLST 和 OASIS$\rightarrow$NLST 则在不同配准任务上测试。加入 SAM 后，前三个 Dice 方向的结果提高，后两个 TRE 方向的误差降低。该实验表明 SAM 在这些测试方向上取得了更好的泛化表现。

\begin{figure}[htbp]
\centering
\includegraphics[width=0.90\textwidth]{ch06/plot_memory_nnn.pdf}\\[-0.2em]
{\thesisfigurefont (a) 回放容量消融实验}\par\vspace{5pt}
\includegraphics[width=0.90\textwidth]{ch06/plot_sam_new.pdf}\\[-0.2em]
{\thesisfigurefont (b) 域内与跨任务泛化实验}
\caption[SAMCL 的回放容量与泛化实验]{SAMCL 的回放容量与泛化实验。（a）不同回放容量下的配准结果，OASIS、CTCT 和 MRCT 报告 Dice，NLST 报告 TRE；（b）采用或不采用锐度感知最小化时的域内与跨任务泛化结果。CTCT 表示腹部 CT--CT 配准，MRCT 表示腹部 MR--CT 配准。}
\label{fig:samcl-ablation-generalization}
\end{figure}

主实验采用固定任务顺序和 LK-UNet 骨干，未系统交换四项任务的到达次序，因此当前结果尚不能反映任务顺序变化下的稳定程度。


\section{持续配准的 MedCL 评测接入与示例设计}
\label{sec:registration-platform-integration}

\subsection{图像对、阶段模型与空间变换接口}

SAMCL 在平台外完成顺序训练和有限回放。MedCL 拟接收一个最终模型、多个阶段模型或对应的配准预测，使用平台侧固定测试图像对与参照评分。用户选择实际完成的任务顺序和阶段；模型方式调用已适配网络，预测方式读取与图像对一一对应的变换结果。训练损失、回放图像对和优化日志不属于必需提交对象。

每个测试样本明确固定图像与移动图像，并记录体素间距、空间方向、预处理及变换约定。本文配准网络采用固定域到移动域的采样映射，以 $I_M\circ\phi$ 得到固定网格上的变形图像；若提交位移场，还需要说明其分量顺序、单位和与恒等映射的关系。平台只能在这些约定一致时执行重采样和评分，不能仅依据张量尺寸推断方向。

\begin{table}[htbp]
\centering
\caption{持续配准接入 MedCL 的评测对象与检查}
\label{tab:samcl-platform-interface}
\thesistablefont
\renewcommand{\arraystretch}{1.16}
\setlength{\tabcolsep}{3pt}
\begin{tabularx}{\textwidth}{@{}p{2.4cm}XX@{}}
\toprule
对象 & 输入与配置 & 评价或检查 \\
\midrule
任务与阶段 & 器官、模态、图像对标识、任务顺序与完成阶段 & 固定任务含义，匹配阶段和测试计划 \\
模型或预测 & 适配结构的模型权重，或样本对应的空间变换 & 输入输出形状、变换方向、空间与单位 \\
结构参照 & 固定和移动图像对应的结构标签 & 对标签执行规定变换，计算结构 Dice \\
地标参照 & 对应地标及物理坐标信息 & 采用匹配方向计算毫米级 TRE \\
结果展示 & 各任务分数、变形图像及有效空间变换 & 配准叠加、网格和阶段变化视图 \\
\bottomrule
\end{tabularx}
\end{table}

\subsection{任务指标与持续学习评价}

OASIS、腹部 CT--CT 和腹部 MR--CT 使用结构 Dice，NLST 使用毫米级 TRE。计算这些指标需提供有效空间变换，或评测允许的变换后结构标签与地标，并与测试参照对应。仅有变形后的影像不足以确定相应结构或地标的位置；仅有变换后标签或地标也不足以评价完整形变场的性质。

各阶段模型在固定测试任务上的结果组成性能矩阵。本章以后向变化表示最终性能与初学性能之差：Dice 的负变化表示重叠精度下降，TRE 的正变化表示地标误差增大。只有最终模型或预测时，报告最终任务结果；具有相应历史阶段结果时，再计算后向变化。缺失结果标为不可计算，不以零值填补，也不跨 Dice 与 TRE 求平均。

\subsection{配准可视化与应用示例}

网页拟并列显示固定图像、原始移动图像和变形图像，通过切片联动、叠加或棋盘格辅助观察对齐关系，并在空间变换有效时显示形变网格。参考标签和地标按数据权限展示，物理单位与变换方向与评分保持一致。不同阶段使用同一测试图像对，有助于观察旧任务的对齐能力如何变化。

结果页面拟分别展示 Dice 和 TRE，注明指标方向、单位及对应任务。平台接入需验证模型或预测读取、坐标约定和评分一致性；形变质量指标需在能够读取有效形变场后单独实现。

\section{方法讨论与研究局限}
\label{sec:registration-discussion}

本章以同一模型适应器官、模态和配准关系的顺序变化。不同任务共享空间变换这一预测形式，但图像相似性、结构参照和形变分布不同，因此历史保持与新任务适应需要结合各任务指标解释。

在固定四任务顺序下，SAMCL 减轻了直接顺序训练在前三个任务上的性能下降。与 MER 比较时，两者在 OASIS 上接近，SAMCL 在腹部 CT--CT 和腹部 MR--CT 上的 Dice 较高，MER 在 NLST 上的 TRE 及其后向变化较低。

泛化比较覆盖两项域内测试和三项跨任务测试。主任务序列同时改变器官、模态、配准关系和损失配置，各因素的独立影响尚未分离；当前实验也未考察其他骨干或不同任务顺序。

SAMCL 保存有限历史原始图像对，其存储与计算需求随回放容量变化。现有评价采用 Dice 和 TRE，尚未系统评价形变折叠、平滑性、可逆性及逆一致性。

\section{本章小结}
\label{sec:registration-summary}

本章提出 SAMCL，将有限历史图像对回放、一阶元持续学习与锐度感知优化结合，用于不同器官和模态的持续配准。四任务实验中，SAMCL 改善了直接顺序训练的历史保持，与 MER 的相对表现则随任务和指标变化。容量消融及域内、跨任务测试进一步表明，回放规模与锐度感知更新均会影响模型表现。

本章给出 MedCL 配准评测接入设计，将固定与移动图像、空间变换、结构或地标参照及阶段结果对应起来，分别报告 Dice、TRE 及相应后向变化。
